\documentclass[10pt,a4paper]{article}
\usepackage{amsmath,amssymb,amsthm,graphicx}
\usepackage{makecell}
% Card-local minimal preamble (journal unit sigma-2025-023-macdonald-identities-weyl-kac-denominator-fo, P2-ahl16 precedent).
%% Verbatim from the pinned source: the author theorem environments and macro block,
%% so every macro and environment of the counted spans is available.
% Omitted and why: the documentclass line (the journal class sigma is not on the host;
%% the assembler supplies the class) and page-geometry/font blocks (they fight the
%% forced xelatex engine).
\usepackage{amsmath,amssymb,amsthm,mathtools}

\renewcommand{\textfraction}{0.01}
\renewcommand{\topfraction}{0.99}
\renewcommand\theadalign{bc}
\renewcommand\theadfont{\bfseries}
\renewcommand\theadgape{\Gape[4pt]}
\renewcommand\cellgape{\Gape[4pt]}
\newtheorem{Theorem}{Theorem}[section]
\newtheorem{Corollary}[Theorem]{Corollary}
\newtheorem{Lemma}[Theorem]{Lemma}
\newtheorem{Proposition}[Theorem]{Proposition}
 { \theoremstyle{definition}
\newtheorem{Definition}[Theorem]{Definition}
\newtheorem{Note}[Theorem]{Note}
\newtheorem{Example}[Theorem]{Example}
\newtheorem{Remark}[Theorem]{Remark} }
\numberwithin{equation}{section}
\newcommand{\ti}{\tilde}
\newcommand{\al}{\alpha}
\newcommand{\om}{\omega}
\newcommand{\ga}{\gamma}
\newcommand{\la}{\lambda}
\newcommand{\eps}{\varepsilon}
\newcommand{\si}{\sigma}
\newcommand{\La}{\Lambda}
\newcommand{\Ga}{\Gamma}
\newcommand{\ov}[1]{\overline{#1}}
\newcommand{\os}[2]{\overset{#1}{#2}}
\newcommand{\scb}[2]{\scalebox{#1}{#2}}
\newcommand{\crossp}{\overset{+}{\leadsto}}
\newcommand{\crossn}{\overset{-}{\leadsto}}
\newcommand{\fh}{\mathfrak{h}}
\newcommand{\fg}{\mathfrak{g}}
%
\newcommand{\cC}{\mathcal{C}}
\newcommand{\cF}{\mathcal{F}}
\newcommand{\cU}{\mathcal{U}}
\newcommand{\cM}{\mathcal{M}}
\newcommand{\cA}{\mathcal{A}}
\newcommand{\cB}{\mathcal{B}}
\def\cch{\mathcal H}
\def\ccsc{\mathcal {SC}}
\def\ccdd{\mathcal {DD}}
\def\ccd{\mathcal {D}}
\def\ccp{\mathcal P}
\def\ccc{\mathcal C}
\DeclareMathOperator{\BGP}{BG}
\DeclareMathOperator{\WLP}{WL}
\DeclareMathOperator{\core}{core}
\DeclareMathOperator{\quot}{quot}
\newcommand{\scal}[2]{\langle #1,#2\rangle}
\newcommand{\scald}[2]{(#1,#2)}
\newcommand{\qu}[1]{\qquad\text{#1}\qquad}
\newcommand{\te}[1]{\text{ #1 }}
\newcommand{\Z}{\mathbb{Z}}
\newcommand{\N}{\mathbb{N}}
\newcommand{\C}{\mathbb{C}}
\newcommand{\R}{\mathbb{R}}
\newcommand{\K}{\mathbb{K}}
\newcommand{\Q}{\mathbb{Q}}
\newcommand{\bB}{\mathbb{B}}
\newcommand{\bA}{\mathbb{A}}
\newcommand{\fB}{\mathfrak{B}}
\newcommand{\wB}{\widehat{B}}
\newcommand{\mk}{\mathrm{k}}
\newcommand{\mK}{\mathrm{K}}
\newcommand{\sLa}{\mathsf{\La}}
\newcommand{\cAJ}{\mathcal{A}_J}
\DeclareMathOperator{\Sp}{Sp}
\DeclareMathOperator{\GL}{GL}
\DeclareMathOperator{\OO}{SO}
\renewcommand{\OE}{\mathrm{O}}
\DeclareMathOperator{\I}{I}
\DeclareMathOperator{\oo}{so}
\renewcommand{\oe}{\mathrm{o}^{\text{even}}}
\newcommand{\func}[1]{\operatorname{#1}}
\newcommand{\lra}{\longrightarrow}
\newcommand{\cev}[1]{\reflectbox{\ensuremath{\vec{\reflectbox{\ensuremath{#1}}}}}}
\newcommand{\cI}{\mathcal{I}}
\newcommand\blfootnote[1]{%
	\begingroup
	\renewcommand\thefootnote{}\footnote{#1}%
	\addtocounter{footnote}{-1}%
	\endgroup
}
\newcommand{\syc}{\mathsf{s_C}}
\newcommand{\lb}{\llbracket}
\newcommand{\rb}{\rrbracket}
\usepackage{cleveref}
\crefname{originaldefinition}{original definition}{original definitions}
\Crefname{originaldefinition}{Original definition}{Original definitions}
\makeatletter
\newlabel{def:Phi}{{18.5}{source}{}{sourcepointer-1}{}}
\newlabel{def:Phi@cref}{{[originaldefinition][0][]18.5}{[1][0][]source}{}{}{}}
\makeatother
\crefname{originaldefinition}{original definition}{original definitions}
\Crefname{originaldefinition}{Original definition}{Original definitions}
\makeatletter
\newlabel{def:psi}{{18.1}{source}{}{sourcepointer-2}{}}
\newlabel{def:psi@cref}{{[originaldefinition][0][]18.1}{[1][0][]source}{}{}{}}
\makeatother
\crefname{originallemma}{original lemma}{original lemmas}
\Crefname{originallemma}{Original lemma}{Original lemmas}
\makeatletter
\newlabel{lem:indices}{{18.2}{source}{}{sourcepointer-3}{}}
\newlabel{lem:indices@cref}{{[originallemma][0][]18.2}{[1][0][]source}{}{}{}}
\makeatother
\crefname{originallemma}{original lemma}{original lemmas}
\Crefname{originallemma}{Original lemma}{Original lemmas}
\makeatletter
\newlabel{lem:DDprimeLittlewood}{{18.7}{source}{}{sourcepointer-4}{}}
\newlabel{lem:DDprimeLittlewood@cref}{{[originallemma][0][]18.7}{[1][0][]source}{}{}{}}
\makeatother
\crefname{originalequation}{original equation}{original equations}
\Crefname{originalequation}{Original equation}{Original equations}
\makeatletter
\newlabel{eq:motddprime}{{6.4}{source}{}{sourcepointer-5}{}}
\newlabel{eq:motddprime@cref}{{[originalequation][0][]6.4}{[1][0][]source}{}{}{}}
\makeatother
\crefname{originalequation}{original equation}{original equations}
\Crefname{originalequation}{Original equation}{Original equations}
\makeatletter
\newlabel{eq:gkdd}{{6.7}{source}{}{sourcepointer-6}{}}
\newlabel{eq:gkdd@cref}{{[originalequation][0][]6.7}{[1][0][]source}{}{}{}}
\makeatother
\title{Doubled distinct partitions of type (d,tr,r): quotient shape and explicit size formulas}
\author{Lecouvey, C\textbackslash{}'\{e\}dric, Wahiche, David}
\date{doi:10.3842/SIGMA.2025.023; CC BY 4.0}
\begin{document}
\maketitle
\setcounter{section}{6}
\setcounter{Theorem}{14}
\setcounter{equation}{7}
\noindent\textit{Editorial note.} All delivered mathematical text is the author's own, copied from the pinned journal source file sources/sigma-2025-023-macdonald-identities-weyl-kac-denominator-fo/pinned-original-source.tex (Macdonald Identities, Weyl-Kac Denominator Formulas and Affine Grassmannian Elements); no statement or proof step was rewritten, reordered or removed.\par\medskip
\noindent\textit{Editorial note.} The counted proof is the author's complete argument: The two mirrored arguments ('using the same arguments', lines 2361-2366 and 2368-2373) are executed inline in compressed form, not replaced by the phrase. The statement's two displayed size formulas (6.8)/(6.9) are dispatched in the proof by named-input pointers ('derived using (6.7) and (DD'3)') - same class as a named external theorem invocation, all declared as external navigation pointers: def:psi (18.1), def:Phi (18.5), lem:indices (18.2), lem:DDprimeLittlewood (18.7), eq:motddprime (6.4), eq:gkdd (6.7). The pinned source carries commented-out author working notes inside the span (lines 2334-2339 and 2386-2389); they are carried verbatim and do not render. The display at line 2378 is an unnumbered equation* carrying label eq:nu0typeb (no printed number exists). The census printed field 20.3 is the published-version coordinate; the card face renders Proposition 6.15 by source-text counter stepping, with the in-span numbered equations rendering as (6.8) and (6.9).\par\medskip
\noindent\textit{Editorial note.} The label prop:ddprime names the counted statement, and the proof environment immediately following it in the source is delivered verbatim as its proof.\par\medskip
\noindent\textit{Editorial note.} The counted statement is carried by the author's own Proposition environment in the article source.\par\medskip
\noindent\textit{Editorial note.} Journal-unit assembly: card-local preamble (author theorem environments and macros verbatim), article class, no bibliography in the counted span (the span carries no \textbackslash{}cite commands).\par\medskip
\noindent\textit{Editorial note.} Revision r002: the r001 card-local preamble was malformed - the lane's preamble extractor dropped five structural lines of the author's \textbackslash{}blfootnote definition (\textbackslash{}begingroup, \textbackslash{}addtocounter, \textbackslash{}endgroup, the closing brace, a bare \% comment), leaving an unbalanced brace and an xelatex failure; the extractor is fixed and the preamble is regenerated verbatim from the pinned source. Additionally \textbackslash{}usepackage\{makecell\} is declared via assembly.extra\_packages because the source preamble redefines the journal's tablehead macros (\textbackslash{}theadalign, \textbackslash{}theadfont, \textbackslash{}theadgape, \textbackslash{}cellgape via \textbackslash{}Gape) that the article class does not define; no card-body text changed.\par\medskip
\noindent\textit{Editorial scope.} the counted statement and its own proof environment only\par\medskip
\noindent\textit{Allowed standard prerequisites (editorial).}\begin{itemize}
\item Definition def:Phi (printed 18.5): the Littlewood decomposition Phi\_n(lambda) = (omega, nu\textasciicircum{}(0),...,nu\textasciicircum{}(n-1)).
\item Definition def:psi (printed 18.1): the binary word (c\_k)\_\{k in Z\} of a partition.
\item Lemma lem:indices (printed 18.2): the index correspondence between zeros of the binary word and boxes.
\item Lemma lem:DDprimeLittlewood (printed 18.7): properties (DD'3) and (DD'4) of doubled distinct partitions under transposition.
\item Equation eq:motddprime (printed (6.4)): the symmetry relation c\_\{-k-2\} = 1 - c\_k of the binary word.
\item Equation eq:gkdd (printed (6.7)): the size identity for the doubled word used in the final derivation.
\end{itemize}\medskip
\ifdefined\hypertarget\hypertarget{sourcepointer-1}{}\fi
\noindent\textit{Source reference.} Original definition 18.5: sources/sigma-2025-023-macdonald-identities-weyl-kac-denominator-fo/pinned-original-source.tex, source lines 2006--2006. Navigation pointer invoked at line(s) 2345: the counted text uses that definition of the same article (the Littlewood decomposition), which is not retained in this card, so the reference is declared as a pointer carrying the number the article itself prints for it (the section and the place of the statement inside that section, derived by the census v3.5 rule for the shared section-scoped theorem counter); the cited definition is an allowed input and none of it is reproduced inside the span.\par\medskip
\ifdefined\hypertarget\hypertarget{sourcepointer-2}{}\fi
\noindent\textit{Source reference.} Original definition 18.1: sources/sigma-2025-023-macdonald-identities-weyl-kac-denominator-fo/pinned-original-source.tex, source lines 1732--1732. Navigation pointer invoked at line(s) 2346: the counted text uses that definition of the same article (the binary-word abacus encoding), which is not retained in this card, so the reference is declared as a pointer carrying the number the article itself prints for it (the section and the place of the statement inside that section, derived by the census v3.5 rule for the shared section-scoped theorem counter); the cited definition is an allowed input and none of it is reproduced inside the span.\par\medskip
\ifdefined\hypertarget\hypertarget{sourcepointer-3}{}\fi
\noindent\textit{Source reference.} Original lemma 18.2: sources/sigma-2025-023-macdonald-identities-weyl-kac-denominator-fo/pinned-original-source.tex, source lines 1829--1829. Navigation pointer invoked at line(s) 2352: the counted text invokes that lemma of the same article (the index correspondence between word positions and boxes) as a named input, which is not retained in this card, so the reference is declared as a pointer carrying the number the article itself prints for it (the section and the place of the statement inside that section, derived by the census v3.5 rule for the shared section-scoped theorem counter); the cited lemma is an allowed input and none of it is reproduced inside the span.\par\medskip
\ifdefined\hypertarget\hypertarget{sourcepointer-4}{}\fi
\noindent\textit{Source reference.} Original lemma 18.7: sources/sigma-2025-023-macdonald-identities-weyl-kac-denominator-fo/pinned-original-source.tex, source lines 2048--2048. Navigation pointer invoked at line(s) 2352: the counted text invokes that lemma of the same article (properties (DD'3)/(DD'4) of doubled distinct partitions) as a named input, which is not retained in this card, so the reference is declared as a pointer carrying the number the article itself prints for it (the section and the place of the statement inside that section, derived by the census v3.5 rule for the shared section-scoped theorem counter); the cited lemma is an allowed input and none of it is reproduced inside the span.\par\medskip
\ifdefined\hypertarget\hypertarget{sourcepointer-5}{}\fi
\noindent\textit{Source reference.} Original equation 6.4: sources/sigma-2025-023-macdonald-identities-weyl-kac-denominator-fo/pinned-original-source.tex, source lines 1985--1985. Navigation pointer invoked at line(s) 2356: the counted text cites that numbered equation of the same article, which is not retained in this card, so the reference is declared as a pointer carrying the number the article itself prints for it (the section and the place of the equation inside that section, derived by the census v3.5 rule for the section-scoped equation counter); the cited equation is an allowed input and none of it is reproduced inside the span.\par\medskip
\ifdefined\hypertarget\hypertarget{sourcepointer-6}{}\fi
\noindent\textit{Source reference.} Original equation 6.7: sources/sigma-2025-023-macdonald-identities-weyl-kac-denominator-fo/pinned-original-source.tex, source lines 2172--2172. Navigation pointer invoked at line(s) 2384: the counted text cites that numbered equation of the same article, which is not retained in this card, so the reference is declared as a pointer carrying the number the article itself prints for it (the section and the place of the equation inside that section, derived by the census v3.5 rule for the section-scoped equation counter); the cited equation is an allowed input and none of it is reproduced inside the span.\par\medskip
\begin{Proposition}
	\label{prop:ddprime} Let $n$ be an integer greater or equal to $2$. Take $%
	\bar{\lambda}\in \ccd$ and set $\Phi _{n}\bigl(\bar{\lambda}\bar{\lambda}^{%
		\mathrm{tr}}\bigr)=\smash{\bigl( \omega ,\nu ^{(0)},\dots ,\nu ^{(n-1)}\bigr)} $ $($see
	Definition $\ref{def:Phi})$. Then $\bar{\lambda}$ belongs to \smash{$\cC_{n}^{d,%
		\mathrm{tr},r}$} if and only if
	\begin{itemize}\itemsep=0pt
		\item $\nu^{(j)}=\varnothing$ for any $0\leq j\leq n-2$,
		
		\item let $m=\max\bigl(\lbrace 1\rbrace\cup\bigl\lbrace\bigl(n+\bar{\lambda}_i\bigr)/n \mid
		\bar{\lambda}_i\equiv 0\pmod{n}\bigr\rbrace\bigr)$. Then $\nu^{(n-1)}$ is the rectangle
		${(m-1)\times m}$ partition.
	\end{itemize}
	
	Moreover, in the odd case, setting $\phi _{2n-1}(\omega )=(m_{0},\dots
	,m_{n-2},-m_{n-2},\dots ,-m_{0},0)$ one gets%
	\begin{equation}
	\bigl\lvert \bar{\lambda}\bigr\rvert =\frac{2n-1}{2}\left(
	m^{2}+\sum_{i=0}^{n-2}m_{i}^{2}\right) +\left( -\frac{2n-1}{2}%
	m+\sum_{i=0}^{n-2}\biggl(i-n+\frac{1}{2}\biggr)m_{i}\right) . \label{eq:gkddnuimpair}
	\end{equation}%
	In the even case, by setting $\phi _{2n}(\omega )=(m_{0},\dots
	,m_{n-2},0,-m_{n-2},\dots ,-m_{0},0)$ one gets%
	\begin{equation}
	\bigl\lvert \bar{\lambda}\bigr\rvert =n\left( m^{2}+\sum_{i=0}^{n-2}m_{i}^{2}\right)
	+\left( -nm+\sum_{i=0}^{n-2}(i-n+1)n_{i}\right) .
	\label{eq:gkddnupairbvtquot}
	\end{equation}%
	%there exists m$\in\N^*$ such that $\nu^{(n-1)}$ is $(m-1)\times m$ partition and for any $j$ $\nu^{(j)}=\varnothing$ fo
	%Set $\bigl(\bar{\lambda}\bar{\lambda}\bigr)^\mathrm{tr}\in\ccdd^{\mathrm{tr}}$ such that $\cch(\bar{\lambda}\bar{\lambda}^\mathrm{tr})$ has only one hook length equal to $n$ if and only if the Littlewood decomposition of $bar{\lambda}\bar{\lambda})^\mathrm{tr}}$ is of the form:
	%\begin{itemize}
	%\item if $n$ is odd, $\Phi_n(\bar{\lambda}\bar{\lambda}^\mathrm{tr}}$
	%\end{itemize}
	%$\#\lbrace h=n,h\in\cch\bigl(\bar{\lambda}\bar{\lambda}\bigr)^\mathrm{tr})\rbrace=1$.
\end{Proposition}

\begin{proof}
	Consider $\bar{\lambda}\in\cC_{n}^{d,\mathrm{tr},r}$, the Littlewood
	decomposition $\Phi_n\bigl(\bar{\lambda}\bar{\lambda}^\mathrm{tr}%
	\bigr)=\bigl(\omega,\nu^{(0)},\dots,\nu^{(n-1)}\bigr)$ (see Definition \ref{def:Phi})
    and $\psi\bigl(\bar{\lambda}\bar{\lambda}^\mathrm{tr}\bigr)=(c_k)_{k\in\Z}$
	its corresponding word (see Definition \ref{def:psi}). Since $n/2\not\in\bar{%
		\lambda}$, the main diagonal $D\bigl(\bar{\lambda}\bar{\lambda}^%
	\mathrm{tr}\bigr)$ of $\bigl(\bar{\lambda}\bar{\lambda}\bigr)^\mathrm{tr}$ has no box of hook length $n$.
	If there exists $i\in\lbrace 0,\dots,n-2\rbrace$ such that $\nu^{(i)}\neq
	\varnothing$, there exists at least one box in $\nu^{(k)}$ whose hook length
	is equal to $1$. By Lemma \ref{lem:indices} and by property $(DD^{\prime }4)$
	from Lemma \ref{lem:DDprimeLittlewood}, there exists $m\in\Z$ such that $%
	c_{mn+i}=1$ and $c_{(m+1)n+i}=0$. Note that the hook length of the box
	$(mn+i,(m+1)n+i)$ is equal to $n$. If $%
	i\not\equiv-i-2\pmod{n}$, that is, if $i\neq n/2$, by \eqref{eq:motddprime},
	one gets that ${c_{-(m+1)n-i-2}=1}$ and ${c_{-mn-i-2}=0}$. Moreover, one of the
	two boxes corresponding to the pairs of indices $(mn+i,(m+1)n+i)$ and $%
	(-(m+1)n-i-2,-mn-i-2)$ is strictly above the main diagonal $D\bigl(\bar{\lambda}\bar{\lambda}^\mathrm{tr}\bigr)$,
	which contradicts the fact that $\smash{\bar{\lambda}\in\cC_{n}^{d,\mathrm{tr},r}}$.
	If $i=n/2$, which implies that $n$ is even in this case, using the same
	arguments, $\smash{\bar{\lambda}\in\cC_{n}^{d,\mathrm{tr},r}}$ implies that there is
	no hook of length $n$ apart maybe on the diagonal~$D\bigl(\bar{\lambda}\bar{\lambda}^\mathrm{tr}\bigr)$. Since that $%
	n/2\not\in \bar{\lambda}$, we have that $\nu^{(n/2)}=\varnothing$. Therefore, \smash{$%
	\bar{\lambda}\in\cC_{n}^{d,\mathrm{tr},r}$} implies that~\smash{${\nu^{(i)}=\varnothing}$}
	for any $0\leq j\leq n-2$.
	
	Now we prove the second part of the statement. Let $d$ be the size of the
	Durfee square of~$\nu^{(n-1)}$. Using the same arguments as above, if $%
	\nu^{(n-1)}$ has a hook length equal to $1$ in any row but the $d+1$-th row,
	the properties of symmetry along the main diagonal and the property~$(DD^{\prime }4)$
    from Lemma \ref{lem:DDprimeLittlewood} imply that $%
	\nu^{(n-1)}$ contains a unique box whose hook length is equal to $1$.
	Therefore, $\nu^{(n-1)}$ is a rectangle of size $m\times(m-1)$ with $%
	m=\max\bigl(\lbrace 1\rbrace\cup\bigl\lbrace\bigl(n+\bar{\lambda}_i\bigr)/n \mid\bar{\lambda%
	}_i\equiv 0\pmod{n}\bigr\rbrace\bigr)$ or equivalently $m=\max(k+1\mid
	c_{kn+n-1}= 0)$. Note that the word to $\nu^{(n-1)}$ is
	\begin{equation*} \label{eq:nu0typeb}
	\psi\bigl(\nu^{(n-1)}\bigr)=\dots 0\underbrace{1\dots 1}_{m-1}\underbrace{0\mid 0\dots
		0}_{m}1\dots.
	\end{equation*}
	
	Equalities \eqref{eq:gkddnuimpair} and \eqref{eq:gkddnupairbvtquot} are
	derived using \eqref{eq:gkdd} and Lemma \ref{lem:DDprimeLittlewood} $%
	(DD^{\prime }3)$.
	% Let $m=\max\left( (\bar{\lambda}_i-n+1)/n,\bar{\lambda}_i\equiv n-1\pmod{n}\right)$ or equivalently $m=\max(k\mid c_{kn+n-1)= 0,k\in\N)$.
	%By Lemma \ref{lem:indices} and by property $(DD'4)$ from Lemma \ref{lem:DDprimeLittlewood}, it implies that there is a box $s'$ in $\nu^{(n-2-k)}$ such that $h_{s'}=1$.
	%If there exists $k\in\lbrace 0,\dots,n-2\rbrace$ such that $\nu^{(k)}\neq \varnothing$, there exists $i,j$ such that $\bigl(\bar{\lambda}\bar{\lambda}\bigr)_i
	%Moreover, by the property $(DD'4)$ from Lemma \ref{lem:DDprimeLittlewood},
\end{proof}

\end{document}
